\documentclass[11pt]{article}
\usepackage[a4paper,margin=27mm]{geometry}
\usepackage{amsmath,amssymb,amsthm,xcolor,booktabs,array}
\newcommand{\PROVED}{\textcolor{green!40!black}{\textbf{PROVED}}}
\newcommand{\OPEN}{\textcolor{red!70!black}{\textbf{OPEN}}}
\newcommand{\AUDIT}{\textcolor{orange!70!black}{\textbf{AUDIT}}}
\newtheorem{theorem}{Theorem}
\newtheorem{corollary}{Corollary}
\title{The Gamma--Dirichlet metric on an arithmetic threshold fiber (v120)}
\date{13 August 2026}

\begin{document}
\maketitle

\section{Pullback of one Gamma resolvent}

Fix $0<\delta<\frac12\log(1+1/N)$ and $E=L^2(0,\delta)$.
For $1\le n\le N$ put $x_n=\alpha-\log n$ and define the translation
isometry
\[
 (T_ne)(x_n+u)=e(u),\qquad 0<u<\delta.
\]
The intervals $x_n+(0,\delta)$ are pairwise disjoint.  For
$v=(v_1,\ldots,v_N)\in\mathbb C^N$ set
\begin{equation}
 \mathfrak S_N(v\otimes e)=\sum_{n\le N}v_nT_ne.             \tag{1}
\end{equation}
This extends to an isometry $\mathbb C^N\widehat\otimes E\to L^2(\mathbb R)$.

For $a>0$ let $K_a$ be convolution by $e^{-a|t|}$.  Its Fourier multiplier
is $2a/(a^2+\tau^2)$.  Define the cell quantities
\begin{align}
 Q_{a,0}(e)&=\int_0^\delta\!\!\int_0^\delta
 e^{-a|u-v|}e(u)\overline{e(v)}\,du\,dv,\notag\\
 M_{a,+}(e)&=\int_0^\delta e^{au}e(u)\,du,\qquad
 M_{a,-}(e)=\int_0^\delta e^{-au}e(u)\,du.                   \tag{2}
\end{align}

\begin{theorem}[Exact Gamma--Dirichlet pullback --- \PROVED]
For every $v\in\mathbb C^N$ and $e\in E$,
\begin{align}
 \langle K_a\mathfrak S_N(v\otimes e),
               \mathfrak S_N(v\otimes e)\rangle
={}&\left(\sum_{n\le N}|v_n|^2\right)Q_{a,0}(e)\notag\\
 &+2\Re\left[
 M_{a,+}(e)\overline{M_{a,-}(e)}
 \sum_{n>m}v_n\overline{v_m}\left(\frac mn\right)^a
 \right].                                                   \tag{3}
\end{align}
Equivalently, every off-diagonal Gamma coupling on the threshold fiber is
the explicit Dirichlet ratio kernel $(m/n)^a$ multiplied by the two
oriented cell moments in (2).
\end{theorem}

\begin{proof}
The diagonal terms in the double sum obtained from (1) give $Q_{a,0}$.
If $n>m$, then $x_n<x_m$ and
\[
 (x_m+v)-(x_n+u)=\log(n/m)+v-u>0
\]
because $\log(n/m)>\delta$.  Therefore
\[
 e^{-a|(x_n+u)-(x_m+v)|}
 =\left(\frac mn\right)^ae^{au}e^{-av}.
\]
The $(n,m)$ term is the product in (3), and the $(m,n)$ term is its complex
conjugate.  Summation proves the formula.
\end{proof}

\section{A positive source Gram for the complete Gamma place}

Put
\begin{equation}
 \mathcal Q_{a,N}:=\frac2aI-\mathfrak S_N^*K_a\mathfrak S_N
 \quad\text{on }\mathbb C^N\widehat\otimes E.                \tag{4}
\end{equation}

\begin{theorem}[Positive Gamma source metric --- \PROVED]
For every $a>0$,
\begin{equation}
                         0\le\mathcal Q_{a,N}\le\frac2aI.    \tag{5}
\end{equation}
For $a_j=2j+\frac12$, the increasing form sum
\begin{equation}
 \mathcal G_{\Gamma,N}^{\rm src}
 :=\sum_{j\ge0}\mathcal Q_{a_j,N}                            \tag{6}
\end{equation}
is closed on its natural domain and satisfies the exact pullback identity
\begin{equation}
 \boxed{
 \mathcal G_{\Gamma,N}^{\rm src}
 =\mathfrak S_N^*G_\Gamma\mathfrak S_N.}                    \tag{7}
\end{equation}
Thus the archimedean part of the threshold fiber has a positive
source-defined Gram before any Schur inversion or sign assumption.
\end{theorem}

\begin{proof}
The Fourier multiplier of $\frac2aI-K_a$ is
\[
 \frac2a-\frac{2a}{a^2+\tau^2}
 =\frac{2\tau^2}{a(a^2+\tau^2)}\ge0.
\]
Compression by the isometry $\mathfrak S_N$ proves (5).  Summing these
multipliers over $a_j$ gives the defining multiplier $g_\Gamma$.
Monotone convergence of positive closed forms proves (6)--(7).
\end{proof}

For the logarithmic incidence vector
\begin{equation}
 c_N(n)=\frac{\Lambda(n)}{\sqrt n},                           \tag{8}
\end{equation}
Theorem 1 becomes a completely explicit cross-place identity:
\begin{align}
 \langle K_a\mathfrak S_N(c_N\otimes e),
                  \mathfrak S_N(c_N\otimes e)\rangle
={}&V_{N,1}Q_{a,0}(e)\notag\\
 &+2\Re\left[
 M_{a,+}(e)\overline{M_{a,-}(e)}
 \sum_{n>m}\frac{\Lambda(n)\Lambda(m)}{\sqrt{mn}}
              \left(\frac mn\right)^a
 \right].                                                   \tag{9}
\end{align}
All terms in (9) are absolutely finite.  No zero of $\xi$ and no desired
row-(d) sign appears.

\section{Relation with the shorted return}

Equations (7)--(9) construct the Gamma part of the source metric
$\mathfrak G_N$ requested in v119 (12).  The Euler part is the finite
divisibility action $\mathcal C_N$ transported by the same isometry, while
the two moment corrections are the terminal graph of v117.  Consequently
the joint Gamma--Euler--Tate source Gram is now explicit as a closed form.

What is not automatic is the direction needed for the inverse metric:
\begin{equation}
 (\mathfrak S_N)^*A_{00}^{-1}\mathfrak S_N
 \le \mathfrak G_N^{\rm GET}.                               \tag{10}
\end{equation}
The equality (7) concerns compression of $G_\Gamma$ to the boundary fiber.
The left side of (10) is the inverse of the shorted complete physical form;
compression and inversion cannot be interchanged.  A proof of (10) must
control the coupling of the boundary fiber to its orthogonal complement.
This is now the only place where the continuum old space remains.

\begin{center}
\begin{tabular}{>{\raggedright\arraybackslash}p{.55\linewidth}
                >{\centering\arraybackslash}p{.15\linewidth}
                >{\raggedright\arraybackslash}p{.20\linewidth}}
\toprule
Statement & Status & Location \\
\midrule
Exact exponential-kernel pullback & \PROVED & (3) \\
Positive one-cell Gamma source Gram & \PROVED & (4)--(5) \\
Complete Gamma pullback to the arithmetic fiber & \PROVED & (6)--(7) \\
Explicit incidence/Gamma coupling & \PROVED & (9) \\
Shorted inverse-metric domination & \OPEN & (10) \\
Condition $G$, row (d) & \OPEN & after joint GET domination \\
\bottomrule
\end{tabular}
\end{center}

\end{document}
